\section{F5 · Kjerneregel, gradient og optimering}

\textbf{Kjernen:} Gradienten peker dit funksjonen stiger raskest og står vinkelrett på nivåkurvene. I et minimum er gradienten null. Gradient descent går i retning $-\nabla f$, og Newton løser $\nabla f=\vb 0$ med Hessematrisen. Med bibetingelse må $\nabla f$ være parallell med $\nabla g$ (Lagrange).

\subsection{Kjerneregelen}
\begin{teorem}[kjerneregelen]
$\vb F(\vb x)=\vb G(\vb U(\vb x))$ med $\vb U:\R^m\to\R^n$ deriverbar i $\vb a$ og $\vb G:\R^n\to\R^p$ deriverbar i $\vb U(\vb a)$. Da er
\[
\underbrace{\vb F'(\vb a)}_{p\times m}=\underbrace{\vb G'(\vb U(\vb a))}_{p\times n}\ \underbrace{\vb U'(\vb a)}_{n\times m}.
\]
\textbf{Med skalarer:} for $f(x,y,z)=g(u(x,y,z),v(x,y,z))$ er $\pd fx=\pd gu\pd ux+\pd gv\pd vx$.
\textbf{Langs en kurve:} $h(t)=f(\vb c(t))$ gir $h'(t)=f'(\vb c(t))\,\vb c'(t)=\ip{\nabla f(\vb c(t))}{\vb c'(t)}$.
\end{teorem}
\begin{fallgruve}[øving 5]
Regn ut $\vb U'$ og $\vb G'$ hver for seg, sett inn $\vb U(\vb a)$ i $\vb G'$ (ikke $\vb a$), og gang så sammen. I et bestemt punkt er dette ofte enklere enn å skrive ut $\vb F$ og derivere direkte.
\end{fallgruve}

\subsection{Optimering: begreper}
\begin{defn}[minimerer]
Problemet er $\vb x^*\in\argmin_{\vb x\in\R^n}f(\vb x)$ for en kostnadsfunksjon (skalarfelt) $f:\R^n\to\R$. For å maksimere $h$ minimerer du $f=-h$. $\vb x^*$ er en \emph{global} minimerer hvis $f(\vb x^*)\le f(\vb x)$ for alle $\vb x$, og en \emph{lokal} minimerer hvis det gjelder i en omegn, for eksempel $\{\vb x:\norm{\vb x-\vb x^*}_2\le R\}$.
\end{defn}

\begin{defn}[gradient, retningsderivert og Hessematrise]
\[
\nabla f(\vb y)=\Bigl(\pd f{x_1}(\vb y),\dots,\pd f{x_n}(\vb y)\Bigr)^\T\in\R^n,\qquad
D_{\vb r}f(\vb y)=\lim_{h\to0}\frac{f(\vb y+h\vb r)-f(\vb y)}{h}\quad(\norm{\vb r}_2=1),
\]
\[
\Hess f(\vb x)=\bigl(\nabla f\bigr)'(\vb x)=\Bigl(\frac{\partial^2f}{\partial x_i\partial x_j}(\vb x)\Bigr)_{i,j=1}^n.
\]
Gradienten er en kolonnevektor, mens Jacobimatrisen $f'(\vb y)=\nabla f(\vb y)^\T$ er en radvektor. For $f\in C^2$ er Hessematrisen symmetrisk.
\end{defn}

\begin{teorem}[gradientens egenskaper]
\begin{enumerate}
  \item Er $f$ deriverbar, så er $D_{\vb r}f(\vb y)=\ip{\nabla f(\vb y)}{\vb r}$.
  \item Med Cauchy--Schwarz: $\nabla f$ peker i retningen med raskest stigning, og $-\nabla f$ i retningen med raskest nedgang. Største stigning er $\norm{\nabla f}_2$.
  \item \textbf{Nivåmengder:} $\{\vb x:f(\vb x)=c\}$ (nivåkurve for $n=2$, nivåflate for $n=3$). Ligger kurven $\vb c(t)$ i en nivåmengde, er $\ip{\nabla f(\vb c(t))}{\vb c'(t)}=0$, så $\nabla f$ står vinkelrett på nivåmengden. Bevis: deriver $f(\vb c(t))=\alpha$ med kjerneregelen.
  \item \textbf{Nødvendig betingelse:} er $\vb y$ et lokalt minimum (eller maksimum) i det indre av definisjonsmengden, så er $\nabla f(\vb y)=\vb 0$ (kritisk eller stasjonært punkt). Det omvendte gjelder ikke: et kritisk punkt kan også være et maksimum eller et sadelpunkt.
\end{enumerate}
\end{teorem}

\begin{formel}[Nyttige gradienter (øving 5)]
\begin{center}
\begin{tabular}{@{}lll@{}}
\toprule
$f(\vb x)$ & $\nabla f(\vb x)$ & $\Hess f(\vb x)$\\
\midrule
$\ip{\vb a}{\vb x}$ & $\vb a$ & $0$\\
$\norm{\vb x}_2$ & $\vb x/\norm{\vb x}_2$ \ ($\vb x\neq\vb 0$) & \\
$\norm{A\vb x}_2^2$ & $2A^\T A\vb x$ & $2A^\T A$\\
$\vb x^\T A\vb x$ & $(A+A^\T)\vb x$ \ ($=2A\vb x$ hvis $A=A^\T$) & $A+A^\T$\\
$\tfrac12\norm{A\vb x-\vb b}_2^2$ & $A^\T(A\vb x-\vb b)$ & $A^\T A$\\
$\tfrac12\norm{A\vb x-\vb y}_2^2+\mu\norm{\vb x}_2^2$ & $A^\T(A\vb x-\vb y)+2\mu\vb x$ & $A^\T A+2\mu I$\\
\bottomrule
\end{tabular}
\end{center}
Den siste er \textbf{Tikhonov-regularisering} (ridge-regresjon). Minimumet løser $(A^\T A+2\mu I)\vb x=A^\T\vb y$, som er entydig løsbart for alle $\mu>0$, også når $A$ ikke har full rang. Når $\mu\to0^+$, går løsningen mot $A^\dagger\vb y$.
\end{formel}

\begin{teorem}[klassifisering (nyttig, men ikke gjennomgått i F5)]
Er $\nabla f(\vb x^*)=\vb 0$ og $\Hess f(\vb x^*)$ positivt definitt (alle egenverdier $>0$), er $\vb x^*$ et lokalt minimum. Negativt definitt gir lokalt maksimum, og egenverdier med begge fortegn gir et sadelpunkt. For $2\times2$: $\det>0$ og $f_{xx}>0$ gir minimum.
\end{teorem}

\subsection{Gradient descent}
\begin{oppskrift}[gradient descent (bratteste nedstigning)]
Gitt $f$, $\nabla f$, startpunkt $\vb x^{(0)}$ og steglengder $s_k$:
\[
\vb x^{(k+1)}=\vb x^{(k)}-s_k\nabla f(\vb x^{(k)}),\qquad k=0,1,2,\dots
\]
\begin{itemize}
  \item \textbf{Steglengde:} konstant $s$ (enkelt). Eksakt linjesøk $\min_t f(\vb x-t\nabla f(\vb x))$ er et nytt optimeringsproblem og som regel for dyrt. I praksis brukes \textbf{backtracking}: prøv $\beta^\ell u$, $\ell=0,1,2,\dots$ (med $0<\beta<1$), og ta det første steget som gir nok nedgang (Armijo-betingelsen, ofte skrevet $f(\vb x-s\nabla f)\le f(\vb x)-c\,s\norm{\nabla f}_2^2$ med for eksempel $c=10^{-4}$).
  \item \textbf{Stopp} når $\norm{\nabla f(\vb x^{(k)})}_2<\varepsilon$ eller etter et maksimalt antall steg.
  \item For stor $s$: iteratene svinger eller divergerer, og $f$ kan øke. For liten $s$: konvergensen blir treg.
  \item Starter du i et kritisk punkt, er $\nabla f=\vb 0$ og metoden står stille (prøveeksamen, oppg.~5c).
\end{itemize}
\end{oppskrift}

\begin{oppskrift}[Newton for optimering]
Bruk Newton (F4) på $\vb F=\nabla f$, slik at Jacobimatrisen blir Hessematrisen:
\[
\Hess f(\vb x^{(k)})\,\Delta^{(k)}=-\nabla f(\vb x^{(k)}),\qquad \vb x^{(k+1)}=\vb x^{(k)}+\Delta^{(k)}.
\]
Metoden trenger andrederiverte og at $\Hess f(\vb x^{(k)})$ er inverterbar. Den konvergerer kvadratisk mot et kritisk punkt, som også kan være et maksimum eller et sadelpunkt. For en kvadratisk $f$ er ett steg nok.
\end{oppskrift}

\begin{oving}
\begin{itemize}[leftmargin=5.6em,labelwidth=5em,labelsep=0.6em,align=left]
\item[\lett] \href{https://github.com/gustavbv/Matte4C---eksamen/blob/main/oppgavesett/F05-gradient-og-optimering/3-hele-oppgaven-i-pensum/007-TMA4135-2016-host-oppg5-oppgave.pdf}{TMA4135 høst 2016, oppg.~5} {\color{graa}(\href{https://github.com/gustavbv/Matte4C---eksamen/blob/main/oppgavesett/F05-gradient-og-optimering/3-hele-oppgaven-i-pensum/007-TMA4135-2016-host-oppg5-svar.pdf}{LF})}: gradient, retningsderivert og nivåkurve
\item[\middels] \href{https://github.com/gustavbv/Matte4C---eksamen/blob/main/oppgavesett/F05-gradient-og-optimering/1-TMA4432-ovinger-og-proveeksamen/002-TMA4432-oving-5-oppg1-oppgave.pdf}{Øving 5, oppg.~1} {\color{graa}(\href{https://github.com/gustavbv/Matte4C---eksamen/blob/main/oppgavesett/F05-gradient-og-optimering/1-TMA4432-ovinger-og-proveeksamen/002-TMA4432-oving-5-oppg1-svar.pdf}{LF})}: kjerneregelen med Jacobimatriser
\item[\middels] \href{https://github.com/gustavbv/Matte4C---eksamen/blob/main/oppgavesett/F05-gradient-og-optimering/1-TMA4432-ovinger-og-proveeksamen/003-TMA4432-oving-5-oppg2-oppgave.pdf}{Øving 5, oppg.~2} {\color{graa}(\href{https://github.com/gustavbv/Matte4C---eksamen/blob/main/oppgavesett/F05-gradient-og-optimering/1-TMA4432-ovinger-og-proveeksamen/003-TMA4432-oving-5-oppg2-svar.pdf}{LF})}: gradienter av normer og kvadratiske former, Tikhonov
\end{itemize}
\end{oving}

\begin{oving}
\begin{itemize}[leftmargin=5.6em,labelwidth=5em,labelsep=0.6em,align=left]
\item[\middels] \href{https://github.com/gustavbv/Matte4C---eksamen/blob/main/oppgavesett/F05-gradient-og-optimering/1-TMA4432-ovinger-og-proveeksamen/001-TMA4432-2026-proveeksamen-oppg5-oppgave.pdf}{Prøveeksamen, oppg.~5} {\color{graa}(\href{https://github.com/gustavbv/Matte4C---eksamen/blob/main/oppgavesett/F05-gradient-og-optimering/1-TMA4432-ovinger-og-proveeksamen/001-TMA4432-2026-proveeksamen-oppg5-svar.pdf}{LF})}: gradient descent for hånd, steglengde og Newton
\item[\middels] \href{https://github.com/gustavbv/Matte4C---eksamen/blob/main/oppgavesett/F05-gradient-og-optimering/1-TMA4432-ovinger-og-proveeksamen/004-TMA4432-oving-5-oppg3-oppgave.pdf}{Øving 5, oppg.~3} {\color{graa}(\href{https://github.com/gustavbv/Matte4C---eksamen/blob/main/oppgavesett/F05-gradient-og-optimering/1-TMA4432-ovinger-og-proveeksamen/004-TMA4432-oving-5-oppg3-svar.pdf}{LF})}: gradient descent og Lagrange
\item[\vanskelig] \href{https://github.com/gustavbv/Matte4C---eksamen/blob/main/oppgavesett/F05-gradient-og-optimering/2-Bergmann-eksamener/005-TMA4215-2022-host-oppg2-oppgave.pdf}{TMA4215 høst 2022, oppg.~2} {\color{graa}(\href{https://github.com/gustavbv/Matte4C---eksamen/blob/main/oppgavesett/F05-gradient-og-optimering/2-Bergmann-eksamener/005-TMA4215-2022-host-oppg2-svar.pdf}{LF})}: gradient descent sammenlignet med normallikningen
\item[\vanskelig] \href{https://github.com/gustavbv/Matte4C---eksamen/blob/main/oppgavesett/F05-gradient-og-optimering/4-delvis-i-pensum/154-TMA4215-2019-kont-oppg3-oppgave.pdf}{TMA4215 kont.\ 2019, oppg.~3} {\color{graa}(\href{https://github.com/gustavbv/Matte4C---eksamen/blob/main/oppgavesett/F05-gradient-og-optimering/4-delvis-i-pensum/154-TMA4215-2019-kont-oppg3-svar.pdf}{LF})}: kritiske punkter og Newton for optimering (delvis pensum)
\end{itemize}
\end{oving}


\subsection{Optimering med bibetingelse: Lagrange-multiplikatorer}
\begin{defn}
Problemet er $\min f(\vb x)$ slik at $g(\vb x)=0$. Punkter med $g(\vb x)=0$ er \emph{tillatte} (feasible). Ved et optimum kan vi ikke kreve $\nabla f=\vb 0$, men ingen tillatt retning (tangent til $\{g=0\}$) kan gi forbedring. Derfor må $\nabla f$ stå vinkelrett på nivåmengden, altså være parallell med $\nabla g$.
\end{defn}
\begin{teorem}[Lagrange]
La $f,g:\R^n\to\R$ ha kontinuerlige gradienter. Hvis $\vb x^*$ er et ekstremalpunkt for $f$ under bibetingelsen $g=0$ og $\nabla g(\vb x^*)\neq\vb 0$, finnes en $\lambda\in\R$ (Lagrange-multiplikator) slik at
\[
\nabla f(\vb x^*)=\lambda\nabla g(\vb x^*).
\]
Med \textbf{Lagrangefunksjonen} $\mathcal L(\vb x,\lambda)=f(\vb x)+\lambda g(\vb x)$ betyr $\nabla\mathcal L=\vb 0$ at $\nabla f+\lambda\nabla g=\vb 0$ og $g(\vb x)=0$. Fortegnet på $\lambda$ er bare en konvensjon. Er $\vb x^*$ også et ubetinget minimum, er $\lambda=0$. Kritiske punkter for $\mathcal L$ er som regel sadelpunkter for $\mathcal L$.
\end{teorem}
\begin{oppskrift}[Lagrange for hånd]
\begin{enumerate}
  \item Sett opp $\nabla f=\lambda\nabla g$ og $g=0$: $n+1$ likninger med $n+1$ ukjente.
  \item Løs, ofte ved å eliminere $\lambda$, for eksempel ved å dele to likninger på hverandre eller dele opp i tilfeller.
  \item Ta også med tillatte punkter der $\nabla g=\vb 0$. Der gjelder ikke teoremet.
  \item Regn ut $f$ i alle kandidatene. Er mengden $\{g=0\}$ lukket og begrenset (for eksempel en sirkel), er den største verdien maksimum og den minste minimum.
\end{enumerate}
\end{oppskrift}
\begin{eksempel}[øving 5, oppg.~3: $f=x^2+4y^2-4x$ på $x^2+y^2=1$]
$\nabla f=(2x-4,\,8y)$ og $\nabla g=(2x,\,2y)$. Likningene $2x-4=2\lambda x$, $8y=2\lambda y$ og $x^2+y^2=1$ gir enten $y=0$, $x=\pm1$ eller $\lambda=4$, $x=-\tfrac23$, $y=\pm\tfrac{\sqrt5}3$. Verdiene er $f(1,0)=-3$ (minimum), $f(-1,0)=5$ og $f(-\tfrac23,\pm\tfrac{\sqrt5}{3})=\tfrac{16}3$ (maksimum). Kontrollregnet.
\end{eksempel}
\textbf{Videre (ikke eksamen):} flere bibetingelser $\vb H(\vb x)=\vb 0$, ulikheter $\vb G(\vb x)\le\vb 0$ (KKT), straffemetoder, utvidet Lagrange og indrepunktsmetoder.

\begin{oving}
\begin{itemize}[leftmargin=5.6em,labelwidth=5em,labelsep=0.6em,align=left]
\item[\lett] \href{https://github.com/gustavbv/Matte4C---eksamen/blob/main/oppgavesett/F04-jacobi-og-newton/3-hele-oppgaven-i-pensum/026-TMA4105-2024-var-oppg1-oppgave.pdf}{TMA4105 vår 2024, oppg.~1} {\color{graa}(\href{https://github.com/gustavbv/Matte4C---eksamen/blob/main/oppgavesett/F04-jacobi-og-newton/3-hele-oppgaven-i-pensum/026-TMA4105-2024-var-oppg1-svar.pdf}{LF})}: minimum med én lineær bibetingelse
\item[\middels] \href{https://github.com/gustavbv/Matte4C---eksamen/blob/main/oppgavesett/F05-gradient-og-optimering/3-hele-oppgaven-i-pensum/056-TMA4105-2017-var-oppg3-oppgave.pdf}{TMA4105 vår 2017, oppg.~3} {\color{graa}(\href{https://github.com/gustavbv/Matte4C---eksamen/blob/main/oppgavesett/F05-gradient-og-optimering/3-hele-oppgaven-i-pensum/056-TMA4105-2017-var-oppg3-svar.pdf}{LF})}: største og minste verdi på en sirkel
\item[\middels] \href{https://github.com/gustavbv/Matte4C---eksamen/blob/main/oppgavesett/F05-gradient-og-optimering/3-hele-oppgaven-i-pensum/028-TMA4105-2025-var-oppg7-oppgave.pdf}{TMA4105 vår 2025, oppg.~7} {\color{graa}(\href{https://github.com/gustavbv/Matte4C---eksamen/blob/main/oppgavesett/F05-gradient-og-optimering/3-hele-oppgaven-i-pensum/028-TMA4105-2025-var-oppg7-svar.pdf}{LF})}: minste avstand fra origo til en kurve
\item[\vanskelig] \href{https://github.com/gustavbv/Matte4C---eksamen/blob/main/oppgavesett/F03-kurver-og-flater/3-hele-oppgaven-i-pensum/066-TMA4105-2005-var-oppg5-oppgave.pdf}{TMA4105 vår 2005, oppg.~5} {\color{graa}(\href{https://github.com/gustavbv/Matte4C---eksamen/blob/main/oppgavesett/F03-kurver-og-flater/3-hele-oppgaven-i-pensum/066-TMA4105-2005-var-oppg5-svar.pdf}{LF})}: høyeste punkt på en skjæringskurve
\end{itemize}
\end{oving}

